\begin{abstract}
The process of learning to program is best facilitated through practical coding exercises complemented by timely, constructive feedback. However, providing detailed feedback on students’ code, especially regarding structural design aspects, can be labour-intensive and time-consuming. This thesis presents the development of AutoTA, an automated tool designed to assist in generating feedback by detecting structural patterns within object-oriented code.

AutoTA uses static code analysis to identify structural patterns and represents them as graphs. By recognising these patterns, the tool evaluates the flexibility, maintainability, and readability of the code, extending feedback beyond functional correctness. AutoTA aims to detect structural patterns automatically, adapt to various programming languages, and deliver concise, educational feedback.

The tool employs abstract syntax trees (ASTs), graph theory, and design pattern recognition, with a Large Language Model (LLM) summarising the feedback for clarity and relevance. 

This thesis demonstrates the feasibility and effectiveness of AutoTA in educational settings, providing a foundation for future research in automated code analysis tools.
\end{abstract}